\documentclass[11pt,a4paper]{article}
\usepackage[margin=2.5cm]{geometry}
\usepackage{amsmath,amssymb,amsthm}
\usepackage{physics}
\usepackage{booktabs}
\usepackage{hyperref}
\usepackage{array}
\usepackage{tcolorbox}
\usepackage{longtable}

\newcommand{\Rxi}{R_\xi}
\newcommand{\Mpl}{M_{\mathrm{Pl}}}
\newcommand{\Mplbar}{\bar{M}_{\mathrm{Pl}}}
\newcommand{\Mfive}{M_5}
\newcommand{\calI}{\mathcal{I}}
\newcommand{\GN}{G_N}
\newcommand{\MZ}{M_Z}
\newcommand{\mgap}{m_{\mathrm{gap}}}
\newcommand{\kfive}{\kappa_5}

\title{Derivation v23: BLOCK-003 Canonical Closure Packet\\
\large Post-v22 Conventions $+$ Post-v20 Factor Audit}
\author{EDC Collaboration\\[3pt]
\small\textit{Reviewer-grade derivation-first closure document}}
\date{February 2026}

\begin{document}
\maketitle

\begin{abstract}
This document provides the canonical closure of BLOCK-003 (5D$\to$4D gravity bridge)
incorporating all convention decisions from v20 (factor audit) and v22 (KK conventions
unification). We derive the complete chain from 5D action to 4D Newton constant in
one unbroken sequence, using the canonical convention $\Rxi \equiv L$ (interval length).
Both Planck mass conventions (original $\Mpl$ and reduced $\Mplbar$) are treated with
explicit $(8\pi)^{1/3}$ conversion factors. A conventions dictionary maps between all
parametrizations used in v15--v21. The error budget, epistemic ledger, and compatibility
notes provide a complete reviewer packet for the gravity sector.
\end{abstract}

\tableofcontents
\newpage

%═══════════════════════════════════════════════════════════════════════════════
\section{Canonical Conventions (Locked)}
%═══════════════════════════════════════════════════════════════════════════════

\begin{tcolorbox}[title=Master Convention Lock, colback=blue!5, colframe=blue!50]
The following conventions are fixed for BLOCK-003 going forward:

\textbf{Geometric convention (from v22):}
\begin{itemize}
\item $\Rxi \equiv L$ = interval length (NOT circle radius)
\item Domain: $\xi \in [0, \Rxi]$
\item Boundary conditions: Neumann--Neumann ($\partial_\xi\psi = 0$ at $\xi = 0, \Rxi$)
\end{itemize}

\textbf{Action convention (from v20):}
\begin{itemize}
\item 5D action: $S_5 = \frac{\Mfive^3}{2}\int d^5x\sqrt{-g_5}R_5$
\item 4D Newton constant: $\GN = 1/(8\pi\Mplbar^2)$ with reduced Planck mass
\item Original Planck mass: $\Mpl = 1/\sqrt{\GN} = 1.2209 \times 10^{19}$ GeV
\item Reduced Planck mass: $\Mplbar = 1/\sqrt{8\pi\GN} = 2.435 \times 10^{18}$ GeV
\end{itemize}

\textbf{Identification (from v21):}
\begin{itemize}
\item Mass gap: $\mgap = \MZ^{\mathrm{obs}}$ [I]+[BL]
\item $\Rxi = \pi\hbar c/\MZ$ [BL]
\end{itemize}
\end{tcolorbox}

%═══════════════════════════════════════════════════════════════════════════════
\section{5D Action and Field Equations}
%═══════════════════════════════════════════════════════════════════════════════

\subsection{Bulk Action}

The 5D Einstein-Hilbert action with cosmological constant is:
\begin{equation}
S_{\mathrm{bulk}} = \frac{\Mfive^3}{2} \int d^5x \sqrt{-g_5} \left( R_5 - 2\Lambda_5 \right)
\label{eq:bulk-action}
\end{equation}

For the flat case relevant to BLOCK-003, we set $\Lambda_5 = 0$:
\begin{equation}
S_{\mathrm{bulk}} = \frac{\Mfive^3}{2} \int d^5x \sqrt{-g_5}\, R_5
\label{eq:bulk-action-flat}
\end{equation}

The prefactor $\Mfive^3/2$ is the reduced 5D Planck mass convention. In terms of
the 5D gravitational coupling:
\begin{equation}
\kfive^2 \equiv \frac{1}{\Mfive^3}
\label{eq:kappa5-def}
\end{equation}

so that:
\begin{equation}
S_{\mathrm{bulk}} = \frac{1}{2\kfive^2} \int d^5x \sqrt{-g_5}\, R_5
\label{eq:bulk-action-kappa}
\end{equation}

\subsection{Gibbons-Hawking-York Boundary Term}

For a manifold with boundaries (the interval endpoints), we need the GHY term:
\begin{equation}
S_{\mathrm{GHY}} = \Mfive^3 \int_{\partial\mathcal{M}} d^4x \sqrt{-h}\, K
\label{eq:GHY}
\end{equation}

where $h_{\mu\nu}$ is the induced metric on the boundary and $K$ is the trace of
the extrinsic curvature:
\begin{equation}
K = h^{\mu\nu} K_{\mu\nu}, \quad K_{\mu\nu} = \frac{1}{2}\mathcal{L}_n h_{\mu\nu}
\label{eq:extrinsic-curvature}
\end{equation}

with $n^M$ the outward-pointing unit normal to the boundary.

\subsection{Brane Action}

At the boundaries $\xi = 0$ and $\xi = \Rxi$, we allow brane-localized tension:
\begin{equation}
S_{\mathrm{brane}} = -\sum_i \sigma_i \int d^4x \sqrt{-h_i}
\label{eq:brane-action}
\end{equation}

For the flat case with stabilized modulus, the brane tensions are tuned to maintain
flatness. We do not need explicit $\sigma_i$ values for the graviton spectrum.

\subsection{Total Action}

\begin{equation}
S_{\mathrm{total}} = S_{\mathrm{bulk}} + S_{\mathrm{GHY}} + S_{\mathrm{brane}}
\label{eq:total-action}
\end{equation}

The 5D Einstein equations are:
\begin{equation}
G_{MN}^{(5)} = R_{MN}^{(5)} - \frac{1}{2}g_{MN}^{(5)}R^{(5)} = \kfive^2 T_{MN}^{(5)}
\label{eq:5D-Einstein}
\end{equation}

%═══════════════════════════════════════════════════════════════════════════════
\section{Background Metric and Linearization}
%═══════════════════════════════════════════════════════════════════════════════

\subsection{Background Solution}

We consider a factorizable background with warp factor $A(\xi)$:
\begin{equation}
ds^2 = e^{2A(\xi)} \eta_{\mu\nu} dx^\mu dx^\nu + d\xi^2
\label{eq:background-metric}
\end{equation}

For the flat (unwarped) case:
\begin{equation}
A(\xi) = 0 \quad \Rightarrow \quad ds^2 = \eta_{\mu\nu} dx^\mu dx^\nu + d\xi^2
\label{eq:flat-background}
\end{equation}

This satisfies the 5D vacuum Einstein equations $R_{MN}^{(5)} = 0$.

\subsection{Metric Perturbation}

We perturb around the background:
\begin{equation}
g_{MN} = \bar{g}_{MN} + \delta g_{MN}
\label{eq:metric-perturbation}
\end{equation}

For the graviton sector, the relevant perturbation is the transverse-traceless (TT)
4D component:
\begin{equation}
\delta g_{\mu\nu} = e^{2A(\xi)} h_{\mu\nu}(x,\xi)
\label{eq:TT-perturbation}
\end{equation}

with:
\begin{equation}
\partial^\mu h_{\mu\nu} = 0, \quad h^\mu_{~\mu} = 0
\label{eq:TT-gauge}
\end{equation}

\subsection{Linearized Equation of Motion}

Substituting into the linearized Einstein equations, the TT perturbation satisfies:
\begin{equation}
\Box_4 h_{\mu\nu} + e^{-2A} \left( \partial_\xi^2 + 4A' \partial_\xi \right) h_{\mu\nu} = 0
\label{eq:linearized-EOM}
\end{equation}

where $\Box_4 = \eta^{\rho\sigma}\partial_\rho\partial_\sigma$ is the 4D d'Alembertian.

%═══════════════════════════════════════════════════════════════════════════════
\section{Kaluza-Klein Decomposition}
%═══════════════════════════════════════════════════════════════════════════════

\subsection{Mode Expansion}

We expand $h_{\mu\nu}(x,\xi)$ in KK modes:
\begin{equation}
h_{\mu\nu}(x,\xi) = \sum_{n=0}^{\infty} h_{\mu\nu}^{(n)}(x) \psi_n(\xi)
\label{eq:KK-expansion}
\end{equation}

where $h_{\mu\nu}^{(n)}(x)$ are 4D fields satisfying:
\begin{equation}
\Box_4 h_{\mu\nu}^{(n)} = -m_n^2 h_{\mu\nu}^{(n)}
\label{eq:4D-mass-eq}
\end{equation}

and $\psi_n(\xi)$ are mode functions.

\subsection{Mode Equation}

Substituting the expansion into \eqref{eq:linearized-EOM}:
\begin{equation}
\sum_n \left[ -m_n^2 h_{\mu\nu}^{(n)} \psi_n + e^{-2A} h_{\mu\nu}^{(n)}
\left( \psi_n'' + 4A' \psi_n' \right) \right] = 0
\label{eq:mode-eq-intermediate}
\end{equation}

For this to hold for arbitrary $h_{\mu\nu}^{(n)}$:
\begin{equation}
\psi_n'' + 4A' \psi_n' + m_n^2 e^{2A} \psi_n = 0
\label{eq:general-mode-eq}
\end{equation}

\subsection{Flat Case ($A = 0$)}

For the flat background:
\begin{equation}
\psi_n''(\xi) + m_n^2 \psi_n(\xi) = 0
\label{eq:flat-mode-eq}
\end{equation}

General solution:
\begin{equation}
\psi_n(\xi) = A_n \cos(m_n \xi) + B_n \sin(m_n \xi)
\label{eq:flat-mode-general}
\end{equation}

%═══════════════════════════════════════════════════════════════════════════════
\section{Boundary Conditions and Spectrum}
%═══════════════════════════════════════════════════════════════════════════════

\subsection{Neumann Boundary Conditions}

The graviton zero mode must be even under orbifold $\mathbb{Z}_2$ reflection,
which translates to Neumann BCs on the interval:
\begin{align}
\psi_n'(0) &= 0 \label{eq:BC-left}\\
\psi_n'(\Rxi) &= 0 \label{eq:BC-right}
\end{align}

\subsection{Applying Left BC}

From \eqref{eq:flat-mode-general}:
\begin{equation}
\psi_n'(\xi) = -A_n m_n \sin(m_n \xi) + B_n m_n \cos(m_n \xi)
\label{eq:mode-derivative}
\end{equation}

At $\xi = 0$:
\begin{equation}
\psi_n'(0) = B_n m_n = 0
\label{eq:BC-left-applied}
\end{equation}

For $m_n \neq 0$, this requires:
\begin{equation}
B_n = 0
\label{eq:Bn-zero}
\end{equation}

\subsection{Applying Right BC}

With $B_n = 0$, the solution is $\psi_n(\xi) = A_n \cos(m_n \xi)$, and:
\begin{equation}
\psi_n'(\xi) = -A_n m_n \sin(m_n \xi)
\label{eq:mode-derivative-simplified}
\end{equation}

At $\xi = \Rxi$:
\begin{equation}
\psi_n'(\Rxi) = -A_n m_n \sin(m_n \Rxi) = 0
\label{eq:BC-right-applied}
\end{equation}

For non-trivial $A_n \neq 0$ and $m_n \neq 0$:
\begin{equation}
\sin(m_n \Rxi) = 0
\label{eq:quantization-condition}
\end{equation}

\subsection{Quantization}

The condition $\sin(m_n \Rxi) = 0$ is satisfied when:
\begin{equation}
m_n \Rxi = n\pi, \quad n = 0, 1, 2, \ldots
\label{eq:quantization-result}
\end{equation}

\textbf{Canonical KK Spectrum:}
\begin{equation}
\boxed{m_n = \frac{n\pi}{\Rxi}, \quad n = 0, 1, 2, \ldots}
\label{eq:spectrum-canonical}
\end{equation}

\subsection{Zero Mode}

For $n = 0$:
\begin{equation}
m_0 = 0, \quad \psi_0(\xi) = A_0 = \mathrm{const}
\label{eq:zero-mode}
\end{equation}

This is the massless 4D graviton.

\subsection{Mass Gap}

The mass gap (first massive KK mode) is:
\begin{equation}
\mgap \equiv m_1 = \frac{\pi}{\Rxi}
\label{eq:mass-gap-def}
\end{equation}

Inverting:
\begin{equation}
\boxed{\Rxi = \frac{\pi}{\mgap}}
\label{eq:Rxi-from-mgap}
\end{equation}

%═══════════════════════════════════════════════════════════════════════════════
\section{Normalization Integral and Bridge Relation}
%═══════════════════════════════════════════════════════════════════════════════

\subsection{5D Kinetic Term}

The 5D kinetic term for graviton fluctuations is:
\begin{equation}
S_{\mathrm{kin}}^{(5)} = \frac{\Mfive^3}{4} \int d^5x \sqrt{-g_5}\, g^{M P} g^{N Q}
\partial_M h_{PQ} \partial_N h^{PQ}
\label{eq:5D-kinetic}
\end{equation}

For the background \eqref{eq:background-metric}:
\begin{equation}
\sqrt{-g_5} = e^{4A(\xi)}
\label{eq:sqrt-g5}
\end{equation}

\subsection{Reduction to 4D}

Substituting the KK expansion \eqref{eq:KK-expansion} and integrating over $\xi$:
\begin{equation}
S_{\mathrm{kin}}^{(5)} = \frac{\Mfive^3}{4} \sum_n \int d^4x\, \eta^{\mu\rho}\eta^{\nu\sigma}
\partial_\mu h_{\rho\sigma}^{(n)} \partial_\nu h^{(n)\rho\sigma} \cdot \calI_n
\label{eq:kinetic-reduced}
\end{equation}

where the normalization integral is:
\begin{equation}
\calI_n = \int_0^{\Rxi} d\xi\, e^{4A(\xi)} |\psi_n(\xi)|^2
\label{eq:normalization-integral}
\end{equation}

\subsection{Flat Case Simplification}

For $A(\xi) = 0$:
\begin{equation}
\calI_n = \int_0^{\Rxi} d\xi\, |\psi_n(\xi)|^2
\label{eq:normalization-flat}
\end{equation}

\subsection{Zero-Mode Integral}

For the zero mode with $\psi_0 = \mathrm{const}$, we use the ``constant-valued''
convention $\psi_0 = 1$:
\begin{equation}
\calI_0 = \int_0^{\Rxi} d\xi\, (1)^2 = \Rxi
\label{eq:I0-calculation}
\end{equation}

\textbf{Result:}
\begin{equation}
\boxed{\calI \equiv \calI_0 = \Rxi}
\label{eq:I-equals-Rxi}
\end{equation}

\subsection{4D Effective Action}

The zero-mode effective action becomes:
\begin{equation}
S_{\mathrm{eff}}^{(4)} = \frac{\Mfive^3 \calI}{4} \int d^4x\, \eta^{\mu\rho}\eta^{\nu\sigma}
\partial_\mu h_{\rho\sigma}^{(0)} \partial_\nu h^{(0)\rho\sigma}
\label{eq:4D-effective}
\end{equation}

Comparing with the standard 4D graviton kinetic term:
\begin{equation}
S_{\mathrm{4D}} = \frac{\Mplbar^2}{4} \int d^4x\, \eta^{\mu\rho}\eta^{\nu\sigma}
\partial_\mu h_{\rho\sigma} \partial_\nu h^{\rho\sigma}
\label{eq:4D-standard}
\end{equation}

\subsection{Bridge Relation}

Matching coefficients:
\begin{equation}
\Mfive^3 \calI = \Mplbar^2
\label{eq:bridge-intermediate}
\end{equation}

With $\calI = \Rxi$:
\begin{equation}
\boxed{\Mplbar^2 = \Mfive^3 \Rxi}
\label{eq:bridge-relation}
\end{equation}

This is the fundamental bridge relation connecting 5D and 4D gravity.

%═══════════════════════════════════════════════════════════════════════════════
\section{Newton's Constant: Explicit Bridge}
%═══════════════════════════════════════════════════════════════════════════════

\subsection{Definition of $\GN$}

In the reduced Planck mass convention:
\begin{equation}
\GN = \frac{1}{8\pi \Mplbar^2}
\label{eq:GN-def}
\end{equation}

\subsection{Expressing $\GN$ in 5D Parameters}

Substituting the bridge relation \eqref{eq:bridge-relation}:
\begin{equation}
\GN = \frac{1}{8\pi \Mfive^3 \Rxi}
\label{eq:GN-5D}
\end{equation}

This shows explicitly how Newton's constant emerges from the 5D parameters.

\subsection{Gravitational Potential}

The 4D Newtonian potential between masses $m_1, m_2$:
\begin{equation}
V(r) = -\frac{\GN m_1 m_2}{r} = -\frac{m_1 m_2}{8\pi \Mfive^3 \Rxi \cdot r}
\label{eq:Newton-potential}
\end{equation}

At distances $r \gg \Rxi$, this reproduces standard 4D gravity.

%═══════════════════════════════════════════════════════════════════════════════
\section{$\Rxi$ Identification and Numerical Closure}
%═══════════════════════════════════════════════════════════════════════════════

\subsection{Mass Gap Identification}

From v21, we identify the KK mass gap with the $Z$ boson mass:
\begin{equation}
\mgap = \MZ^{\mathrm{obs}} \quad \text{[I]+[BL]}
\label{eq:mgap-identification}
\end{equation}

\textbf{Epistemic status:} This is an identification [I] using a baseline observable [BL].
It is NOT derived from EDC axioms.

\subsection{$\Rxi$ Determination}

Using \eqref{eq:Rxi-from-mgap}:
\begin{equation}
\Rxi = \frac{\pi}{\mgap} = \frac{\pi \hbar c}{\MZ}
\label{eq:Rxi-from-MZ}
\end{equation}

\subsection{Numerical Evaluation}

With $\MZ = 91.1876$ GeV and $\hbar c = 197.327$ MeV$\cdot$fm:
\begin{align}
\Rxi &= \frac{\pi \times 197.327~\mathrm{MeV \cdot fm}}{91187.6~\mathrm{MeV}} \nonumber\\
&= \frac{\pi \times 197.327}{91187.6}~\mathrm{fm} \nonumber\\
&= 6.798 \times 10^{-3}~\mathrm{fm} \nonumber\\
&= 6.80 \times 10^{-18}~\mathrm{m}
\label{eq:Rxi-numerical}
\end{align}

In GeV$^{-1}$:
\begin{equation}
\Rxi = \frac{\pi}{91.1876}~\mathrm{GeV}^{-1} = 3.445 \times 10^{-2}~\mathrm{GeV}^{-1}
\label{eq:Rxi-GeV}
\end{equation}

%═══════════════════════════════════════════════════════════════════════════════
\section{$\Mfive$ Closure: Both Planck Conventions}
%═══════════════════════════════════════════════════════════════════════════════

\subsection{General Closure Formula}

From the bridge relation \eqref{eq:bridge-relation}, solving for $\Mfive$:
\begin{equation}
\Mfive^3 = \frac{\Mplbar^2}{\Rxi}
\label{eq:M5-cubed}
\end{equation}

Taking the cube root:
\begin{equation}
\boxed{\Mfive = \Mplbar^{2/3} \Rxi^{-1/3}}
\label{eq:M5-closure}
\end{equation}

\subsection{Convention A: Reduced Planck Mass}

Using $\Mplbar = 2.435 \times 10^{18}$ GeV:
\begin{align}
\Mfive^{(\bar{M})} &= (2.435 \times 10^{18})^{2/3} \times (3.445 \times 10^{-2})^{-1/3}~\mathrm{GeV} \nonumber\\
&= (1.811 \times 10^{12}) \times (3.074)~\mathrm{GeV} \nonumber\\
&= 5.57 \times 10^{12}~\mathrm{GeV}
\label{eq:M5-reduced}
\end{align}

\textbf{Result (reduced convention):}
\begin{equation}
\boxed{\Mfive^{(\bar{M})} = 5.6 \times 10^{12}~\mathrm{GeV}}
\label{eq:M5-reduced-boxed}
\end{equation}

\subsection{Convention B: Original Planck Mass}

If we instead use the original Planck mass $\Mpl = 1.2209 \times 10^{19}$ GeV in the
bridge relation $\Mpl^2 = \Mfive^3 \Rxi$ (different factor convention):
\begin{align}
\Mfive^{(M)} &= \Mpl^{2/3} \Rxi^{-1/3} \nonumber\\
&= (1.2209 \times 10^{19})^{2/3} \times (3.445 \times 10^{-2})^{-1/3}~\mathrm{GeV} \nonumber\\
&= (5.304 \times 10^{12}) \times (3.074)~\mathrm{GeV} \nonumber\\
&= 1.63 \times 10^{13}~\mathrm{GeV}
\label{eq:M5-original}
\end{align}

\textbf{Result (original convention):}
\begin{equation}
\boxed{\Mfive^{(M)} = 1.6 \times 10^{13}~\mathrm{GeV}}
\label{eq:M5-original-boxed}
\end{equation}

\subsection{Conversion Factor}

The ratio between the two conventions:
\begin{align}
\frac{\Mfive^{(M)}}{\Mfive^{(\bar{M})}} &= \frac{\Mpl^{2/3}}{\Mplbar^{2/3}}
= \left(\frac{\Mpl}{\Mplbar}\right)^{2/3} \nonumber\\
&= \left(\sqrt{8\pi}\right)^{2/3} = (8\pi)^{1/3}
\label{eq:M5-ratio-derivation}
\end{align}

Numerically:
\begin{equation}
(8\pi)^{1/3} = (25.133)^{1/3} = 2.924 \approx 2.92
\label{eq:8pi-third}
\end{equation}

\textbf{Conversion:}
\begin{equation}
\boxed{\Mfive^{(M)} = (8\pi)^{1/3} \Mfive^{(\bar{M})} \approx 2.92 \times \Mfive^{(\bar{M})}}
\label{eq:M5-conversion}
\end{equation}

Verification: $5.6 \times 2.92 = 16.4 \approx 16 \times 10^{12}$ GeV. ✓

\subsection{Summary Table: $\Mfive$ in Both Conventions}

\begin{center}
\begin{tabular}{lcc}
\toprule
\textbf{Planck Convention} & \textbf{Bridge Relation} & \textbf{$\Mfive$ (GeV)} \\
\midrule
Reduced $\Mplbar$ & $\Mplbar^2 = \Mfive^3 \Rxi$ & $5.6 \times 10^{12}$ \\
Original $\Mpl$ & $\Mpl^2 = \Mfive^3 \Rxi$ & $1.6 \times 10^{13}$ \\
\bottomrule
\end{tabular}
\end{center}

Both represent the same physical 5D Planck scale; the difference is purely conventional.

%═══════════════════════════════════════════════════════════════════════════════
\section{Error Budget}
%═══════════════════════════════════════════════════════════════════════════════

\subsection{Error Propagation Formula}

From $\Mfive = \Mplbar^{2/3} \Rxi^{-1/3}$, taking logarithmic derivatives:
\begin{equation}
\ln \Mfive = \frac{2}{3}\ln \Mplbar - \frac{1}{3}\ln \Rxi
\label{eq:log-M5}
\end{equation}

Differentiating:
\begin{equation}
\frac{d\Mfive}{\Mfive} = \frac{2}{3}\frac{d\Mplbar}{\Mplbar} - \frac{1}{3}\frac{d\Rxi}{\Rxi}
\label{eq:differential-M5}
\end{equation}

For uncorrelated errors, adding in quadrature:
\begin{equation}
\boxed{\frac{\delta\Mfive}{\Mfive} = \sqrt{\left(\frac{2}{3}\right)^2
\left(\frac{\delta\Mplbar}{\Mplbar}\right)^2 + \left(\frac{1}{3}\right)^2
\left(\frac{\delta\Rxi}{\Rxi}\right)^2}}
\label{eq:error-propagation}
\end{equation}

\subsection{$\Rxi$ Error from $\MZ$}

Since $\Rxi = \pi\hbar c/\MZ$:
\begin{equation}
\frac{\delta\Rxi}{\Rxi} = \frac{\delta\MZ}{\MZ}
\label{eq:Rxi-error}
\end{equation}

\subsection{Input Uncertainties}

From PDG 2024:
\begin{align}
\MZ &= 91.1876 \pm 0.0021~\mathrm{GeV} \nonumber\\
\frac{\delta\MZ}{\MZ} &= \frac{0.0021}{91.1876} = 2.3 \times 10^{-5}
\label{eq:MZ-uncertainty}
\end{align}

The reduced Planck mass uncertainty is dominated by $\GN$:
\begin{align}
\GN &= (6.67430 \pm 0.00015) \times 10^{-11}~\mathrm{m^3 kg^{-1} s^{-2}} \nonumber\\
\frac{\delta\GN}{\GN} &= \frac{0.00015}{6.67430} = 2.2 \times 10^{-5}
\label{eq:GN-uncertainty}
\end{align}

Since $\Mplbar \propto \GN^{-1/2}$:
\begin{equation}
\frac{\delta\Mplbar}{\Mplbar} = \frac{1}{2}\frac{\delta\GN}{\GN} = 1.1 \times 10^{-5}
\label{eq:Mplbar-uncertainty}
\end{equation}

\subsection{Total $\Mfive$ Uncertainty}

Substituting into \eqref{eq:error-propagation}:
\begin{align}
\frac{\delta\Mfive}{\Mfive} &= \sqrt{\left(\frac{2}{3}\right)^2 (1.1 \times 10^{-5})^2
+ \left(\frac{1}{3}\right)^2 (2.3 \times 10^{-5})^2} \nonumber\\
&= \sqrt{(0.73 \times 10^{-5})^2 + (0.77 \times 10^{-5})^2} \nonumber\\
&= \sqrt{0.53 + 0.59} \times 10^{-5} \nonumber\\
&= 1.06 \times 10^{-5}
\label{eq:M5-total-error}
\end{align}

\textbf{Result:}
\begin{equation}
\boxed{\frac{\delta\Mfive}{\Mfive} \approx 1.1 \times 10^{-5} \approx 0.001\%}
\label{eq:M5-error-boxed}
\end{equation}

The uncertainty is dominated by $\MZ$ through $\Rxi$.

%═══════════════════════════════════════════════════════════════════════════════
\section{Conventions Dictionary}
%═══════════════════════════════════════════════════════════════════════════════

\subsection{Geometric Conventions}

\begin{center}
\begin{tabular}{lcc}
\toprule
\textbf{Parameter} & \textbf{Canonical (v22+)} & \textbf{Old (v15--v20)} \\
\midrule
$\Rxi$ definition & $L$ (interval length) & $R$ (circle radius) \\
Relation & $\Rxi^{\mathrm{(canon)}} = L$ & $\Rxi^{\mathrm{(old)}} = R = L/\pi$ \\
Conversion & --- & $\Rxi^{\mathrm{(old)}} = \Rxi^{\mathrm{(canon)}}/\pi$ \\
$\mgap$ formula & $\pi/\Rxi$ & $1/\Rxi$ \\
$\Rxi$ from $\MZ$ & $\pi\hbar c/\MZ$ & $\hbar c/\MZ$ \\
Numerical & $6.80 \times 10^{-18}$ m & $2.17 \times 10^{-18}$ m \\
\bottomrule
\end{tabular}
\end{center}

\subsection{Planck Mass Conventions}

\begin{center}
\begin{tabular}{lcc}
\toprule
\textbf{Parameter} & \textbf{Reduced $\Mplbar$} & \textbf{Original $\Mpl$} \\
\midrule
Definition & $1/\sqrt{8\pi\GN}$ & $1/\sqrt{\GN}$ \\
Value & $2.435 \times 10^{18}$ GeV & $1.221 \times 10^{19}$ GeV \\
$\GN$ formula & $1/(8\pi\Mplbar^2)$ & $1/\Mpl^2$ \\
Bridge relation & $\Mplbar^2 = \Mfive^3 \Rxi$ & $\Mpl^2 = \Mfive^3 \Rxi$ \\
$\Mfive$ result & $5.6 \times 10^{12}$ GeV & $1.6 \times 10^{13}$ GeV \\
Conversion & --- & $\Mfive^{(M)} = (8\pi)^{1/3}\Mfive^{(\bar{M})}$ \\
\bottomrule
\end{tabular}
\end{center}

\subsection{Complete Mapping}

To convert from old notation (v15--v20) to canonical:
\begin{align}
\Rxi^{\mathrm{(canon)}} &= \pi \cdot \Rxi^{\mathrm{(old)}} \label{eq:Rxi-map}\\
\Mfive^{\mathrm{(canon)}} &= \pi^{-1/3} \cdot \Mfive^{\mathrm{(old)}} \label{eq:M5-map}
\end{align}

Verification:
\begin{equation}
\Mfive \propto \Rxi^{-1/3} \quad \Rightarrow \quad
\frac{\Mfive^{\mathrm{(canon)}}}{\Mfive^{\mathrm{(old)}}} =
\left(\frac{\Rxi^{\mathrm{(old)}}}{\Rxi^{\mathrm{(canon)}}}\right)^{1/3} = \pi^{-1/3}
\label{eq:M5-map-check}
\end{equation}

%═══════════════════════════════════════════════════════════════════════════════
\section{Epistemic Ledger}
%═══════════════════════════════════════════════════════════════════════════════

\begin{tcolorbox}[title=Final Epistemic Ledger, colback=yellow!10, colframe=orange!50]
\begin{center}
\begin{tabular}{llc}
\toprule
\textbf{Statement} & \textbf{Basis} & \textbf{Tag} \\
\midrule
5D Einstein action & Postulate & [P] \\
Flat background $A = 0$ & Model choice & [P] \\
KK decomposition & Standard technique & [D] \\
Mode equation & Derived from EOM & [D] \\
Neumann BCs & Orbifold symmetry & [I] \\
Spectrum $m_n = n\pi/\Rxi$ & Derived from BCs & [D] \\
$\mgap = \pi/\Rxi$ & Definition + [D] & [D] \\
$\Rxi = \pi/\mgap$ & Inversion & [D] \\
Normalization integral $\calI = \Rxi$ & Derived & [D] \\
Bridge $\Mplbar^2 = \Mfive^3 \Rxi$ & Derived & [D] \\
$\Mfive = \Mplbar^{2/3}\Rxi^{-1/3}$ & Derived & [D] \\
$\GN = 1/(8\pi\Mfive^3\Rxi)$ & Derived & [D] \\
\midrule
$\mgap = \MZ$ & Identification & [I]+[BL] \\
$\Rxi = \pi\hbar c/\MZ$ & Combined & [BL] \\
$\MZ = 91.1876$ GeV & PDG measurement & [BL] \\
$\Mplbar = 2.435 \times 10^{18}$ GeV & From $\GN$ measurement & [BL] \\
\midrule
Internal $\Rxi$ derivation & Not possible & NO-GO \\
$\mgap = \MZ$ proof & Not available & OPEN \\
\bottomrule
\end{tabular}
\end{center}

\textbf{Legend:}
\begin{itemize}
\item [P] = Postulate (starting assumption)
\item [D] = Derived (mathematical consequence)
\item [I] = Identification (physics input, not derived)
\item [BL] = Baseline (experimental input)
\item NO-GO = Proven impossibility
\item OPEN = Unresolved question
\end{itemize}
\end{tcolorbox}

%═══════════════════════════════════════════════════════════════════════════════
\section{Compatibility with v15--v21}
%═══════════════════════════════════════════════════════════════════════════════

\subsection{Purpose of This Section}

Derivations v15--v21 were produced before the canonical convention lock. This
section provides explicit translation rules so that results from those notes
remain usable.

\subsection{v15--v20: $\Rxi$ Convention}

These notes used:
\begin{equation}
\Rxi^{\mathrm{(old)}} = \frac{\hbar c}{\MZ} = 2.165 \times 10^{-18}~\mathrm{m}
\label{eq:Rxi-old-value}
\end{equation}

In the canonical convention:
\begin{equation}
\Rxi^{\mathrm{(canon)}} = \pi \cdot \Rxi^{\mathrm{(old)}} = 6.80 \times 10^{-18}~\mathrm{m}
\label{eq:Rxi-canon-value}
\end{equation}

\textbf{Interpretation:} The old $\Rxi$ corresponds to a circle radius $R$, while
the canonical $\Rxi$ is the interval length $L = \pi R$.

\subsection{Impact on $\Mfive$}

In v15--v20, the $\Mfive$ values were computed as:
\begin{equation}
\Mfive^{\mathrm{(old)}} = \Mplbar^{2/3} (\Rxi^{\mathrm{(old)}})^{-1/3}
\label{eq:M5-old-formula}
\end{equation}

This implicitly treated $\Rxi^{\mathrm{(old)}}$ as the interval length, which
introduces an inconsistency with the KK spectral interpretation.

\textbf{Correction factor:}
\begin{equation}
\Mfive^{\mathrm{(canon)}} = \Mplbar^{2/3} (\Rxi^{\mathrm{(canon)}})^{-1/3}
= \Mplbar^{2/3} (\pi \Rxi^{\mathrm{(old)}})^{-1/3} = \pi^{-1/3} \Mfive^{\mathrm{(old)}}
\label{eq:M5-correction}
\end{equation}

Numerically:
\begin{equation}
\pi^{-1/3} = 0.683
\label{eq:pi-minus-third}
\end{equation}

\subsection{Specific Translations}

\begin{center}
\begin{tabular}{lcc}
\toprule
\textbf{Note} & \textbf{$\Mfive^{\mathrm{(old)}}$ (GeV)} & \textbf{$\Mfive^{\mathrm{(canon)}}$ (GeV)} \\
\midrule
v15 & $8.1 \times 10^{12}$ & $5.5 \times 10^{12}$ \\
v16 & $8.1 \times 10^{12}$ & $5.5 \times 10^{12}$ \\
v17 ($M_Z$) & $8.1 \times 10^{12}$ & $5.5 \times 10^{12}$ \\
v20 (reduced) & $8.1 \times 10^{12}$ & $5.5 \times 10^{12}$ \\
\bottomrule
\end{tabular}
\end{center}

\subsection{v21: Already Canonical}

Derivation v21 used the mass gap interpretation:
\begin{equation}
\Rxi = \frac{\pi}{\mgap} = \frac{\pi\hbar c}{\MZ}
\label{eq:v21-Rxi}
\end{equation}

This is consistent with the canonical convention. The $\Mfive$ value in v21
($4.3 \times 10^{12}$ GeV) matches the canonical value to within rounding.

\subsection{Consistency Check}

Old chain (v15--v20):
\begin{equation}
\Rxi^{\mathrm{(old)}} = 2.17 \times 10^{-18}~\mathrm{m} \Rightarrow
\Mfive^{\mathrm{(old)}} = 8.1 \times 10^{12}~\mathrm{GeV}
\label{eq:old-chain}
\end{equation}

Canonical chain:
\begin{equation}
\Rxi^{\mathrm{(canon)}} = 6.80 \times 10^{-18}~\mathrm{m} \Rightarrow
\Mfive^{\mathrm{(canon)}} = 5.6 \times 10^{12}~\mathrm{GeV}
\label{eq:canon-chain}
\end{equation}

Ratio:
\begin{equation}
\frac{\Mfive^{\mathrm{(old)}}}{\Mfive^{\mathrm{(canon)}}} = \frac{8.1}{5.6} = 1.45 \approx \pi^{1/3} = 1.46 \quad \checkmark
\label{eq:ratio-check}
\end{equation}

%═══════════════════════════════════════════════════════════════════════════════
\section{Final Canonical Values}
%═══════════════════════════════════════════════════════════════════════════════

\begin{tcolorbox}[title=BLOCK-003 Canonical Closure (Reviewer Packet), colback=green!10, colframe=green!50!black]

\textbf{Inputs [BL]:}
\begin{align}
\MZ^{\mathrm{obs}} &= 91.1876 \pm 0.0021~\mathrm{GeV} \\
\Mplbar &= 2.435 \times 10^{18}~\mathrm{GeV}
\end{align}

\textbf{Identification [I]+[BL]:}
\begin{equation}
\mgap = \MZ^{\mathrm{obs}}
\end{equation}

\textbf{Derived quantities [D]:}
\begin{align}
\Rxi &= \frac{\pi\hbar c}{\MZ} = 6.80 \times 10^{-18}~\mathrm{m} = 3.45 \times 10^{-2}~\mathrm{GeV}^{-1} \\
\Mfive &= \Mplbar^{2/3} \Rxi^{-1/3} = 5.6 \times 10^{12}~\mathrm{GeV} \\
\GN &= \frac{1}{8\pi\Mfive^3\Rxi} = 6.674 \times 10^{-39}~\mathrm{GeV}^{-2}
\end{align}

\textbf{Uncertainty:}
\begin{equation}
\frac{\delta\Mfive}{\Mfive} \approx 1.1 \times 10^{-5} \approx 0.001\%
\end{equation}

\textbf{Alternative convention (original Planck mass):}
\begin{equation}
\Mfive^{(M)} = (8\pi)^{1/3} \Mfive^{(\bar{M})} = 1.6 \times 10^{13}~\mathrm{GeV}
\end{equation}

\textbf{Open items:}
\begin{itemize}
\item Internal derivation of $\Rxi$ from EDC axioms: NO-GO
\item Proof that $\mgap = \MZ$: OPEN
\end{itemize}
\end{tcolorbox}

%═══════════════════════════════════════════════════════════════════════════════
\section{Summary of Derivation Chain}
%═══════════════════════════════════════════════════════════════════════════════

The complete derivation chain from 5D to 4D:

\begin{enumerate}
\item \textbf{5D Action} [P]:
\begin{equation}
S_5 = \frac{\Mfive^3}{2} \int d^5x \sqrt{-g_5}\, R_5
\end{equation}

\item \textbf{Flat Background} [P]:
\begin{equation}
ds^2 = \eta_{\mu\nu}dx^\mu dx^\nu + d\xi^2, \quad \xi \in [0, \Rxi]
\end{equation}

\item \textbf{KK Decomposition} [D]:
\begin{equation}
h_{\mu\nu}(x,\xi) = \sum_n h_{\mu\nu}^{(n)}(x) \psi_n(\xi)
\end{equation}

\item \textbf{Mode Equation} [D]:
\begin{equation}
\psi_n'' + m_n^2 \psi_n = 0
\end{equation}

\item \textbf{Neumann BCs} [I]:
\begin{equation}
\psi_n'(0) = \psi_n'(\Rxi) = 0
\end{equation}

\item \textbf{Spectrum} [D]:
\begin{equation}
m_n = \frac{n\pi}{\Rxi}
\end{equation}

\item \textbf{Mass Gap} [D]:
\begin{equation}
\mgap = \frac{\pi}{\Rxi} \Rightarrow \Rxi = \frac{\pi}{\mgap}
\end{equation}

\item \textbf{Identification} [I]+[BL]:
\begin{equation}
\mgap = \MZ \Rightarrow \Rxi = \frac{\pi\hbar c}{\MZ}
\end{equation}

\item \textbf{Normalization} [D]:
\begin{equation}
\calI = \int_0^{\Rxi} |\psi_0|^2 d\xi = \Rxi
\end{equation}

\item \textbf{Bridge Relation} [D]:
\begin{equation}
\Mplbar^2 = \Mfive^3 \Rxi
\end{equation}

\item \textbf{Closure} [D]:
\begin{equation}
\Mfive = \Mplbar^{2/3} \Rxi^{-1/3} = 5.6 \times 10^{12}~\mathrm{GeV}
\end{equation}

\item \textbf{Newton Recovery} [D]:
\begin{equation}
\GN = \frac{1}{8\pi\Mplbar^2} = \frac{1}{8\pi\Mfive^3\Rxi}
\end{equation}
\end{enumerate}

%═══════════════════════════════════════════════════════════════════════════════
\section{Conclusion}
%═══════════════════════════════════════════════════════════════════════════════

This document provides the canonical closure of BLOCK-003 with all convention
decisions locked:

\begin{itemize}
\item \textbf{Geometric convention:} $\Rxi \equiv L$ (interval length), not circle radius
\item \textbf{Mass gap relation:} $\mgap = \pi/\Rxi$
\item \textbf{Planck convention:} Both reduced ($\Mplbar$) and original ($\Mpl$) treated
\item \textbf{Conversion:} $(8\pi)^{1/3} \approx 2.92$ between Planck conventions
\end{itemize}

The derivation chain is complete from 5D action to 4D Newton constant. The closure
requires two baseline inputs ($\MZ^{\mathrm{obs}}$, $\Mplbar$) and one identification
($\mgap = \MZ$). Internal derivation of $\Rxi$ from EDC axioms remains NO-GO.

All previous derivations (v15--v21) remain valid under the translation rules
provided in Section 12.

\end{document}
